
\section{Model}
\label{sec:model}

Consider a general task that includes an encoding of observed data $x$ into latent variables $z$ through the conditional likelihood $q(z|x)$ (an encoder). Further assume that we observe a variable $c$ which exhibits statistical dependence with $x$ (possibly non-linear dependence). We would like to find a $q(z|x)$ that minimizes our loss function $\mathcal{L}$ from the original task, but that also produces a $z$ independent of $c$.

This is clearly a difficult optimization; independence is a very strong condition. A natural relaxation of this is the minimization of the mutual information $I(z,c)$. We can write our relaxed objective as
\begin{align}
\min_q \mathcal{L}(q,x) + \lambda I(z,c) \label{eq:gen_loss}
\end{align}
where $\lambda$ is a trade-off parameter between the two objectives. $\mathcal{L}$ might involve other variables as well, e.g. labels $y$. Without details of $\mathcal{L}$ and its associated task, we can still provide insight into $I(z,c)$. 

Before continuing it is important to note that all entropic quantities related to $z$ are from the encoding distribution $q$ unless explicitly stated otherwise. In some cases entropies depend on prior distributions, $p(z)$, and this will be explicitly noted.

From properties of mutual information, we have that $I(z,c) = I(z, x) - I(z, x | c) + I(z, c |x)$. Here, we note that $q(z|x)$ is the function that we are optimizing over, and thus the distribution of $z$ solely depends on $x$. Thus,
\begin{align}
I(z, c |x) & = H(z|x) - H(z|x,c) = H(z|x) - H(z|x) = 0.
\end{align}
Using Mutual Information properties and a variational inequality, we can then write the following:
\begin{align}
I(z,c) & = I(z, x) - I(z, x | c)\label{eq:fork_in_road}\\
& = I(z, x) - H(x|c) + H(x | z, c) \\
& \leq I(z, x) - H(x|c) - \mathbb{E}_{x,c,z \sim q}[ \log p(x|z,c)] \label{eq:breakdown} \\
& = \mathbb{E}_{z,x}[ \log q(z|x) - \log q(z)] - H(x|c) - \mathbb{E}_{x,c,z\sim q}[ \log p(x|z,c)]\\
%& = \mathbb{E}_{x}\mathbb{E}_{z\sim q} [ \log q(z | x) - \log q(z) ] - H(x|c) - \mathbb{E}_{x,c,z\sim q}[ \log p(x|z,c)]\\
& = \mathbb{E}_{x}[~KL[~q(z|x)~\|~q(z)~]~  ] - H(x|c) - \mathbb{E}_{x,c,z\sim q}[ \log p(x|z,c)]. \label{eq:penalty}
\end{align}
$H(x|c)$ is a constant and can be ignored. In Eq.~\ref{eq:breakdown} we introduce the variational distribution $p(x|z,c)$ which will play the traditional role of the decoder. 
$I(z,c)$ is thus bounded up to a constant by a divergence and a reconstruction error.

The result is similar in appearance to the bound from Variational Auto-Encoders \cite{kingma2013auto}, wherein we balance the divergence between $q(z|x)$ and a prior $p(z)$ against the reconstruction error. Here our penalty on $I(z,c)$ amounts to encouraging $q(z|x)$ to be close to its marginal $q(z)$, i.e. to vary less across inputs $x$, no matter the form of $q(z|x)$ or $q(z)$.
%We could also interpret this in a coding context using the $I(z, x)$ term from Eq \ref{eq:breakdown}
From a coding viewpoint our penalty encourages the compression of $x$ out of $z$ using the $I(z, x)$ term from Eq \ref{eq:breakdown}.

In both interpretations, these penalties are tempered by conditional reconstruction error. This provides additional intuition; by adding a copy of $c$ into the reconstruction, we ensure that compressing away information in $z$ about $c$ is \emph{not} penalized. In other words, conditional reconstruction combined with compressing regularization leads to invariance w.r.t. to the conditional input.


%The divergence term is minimal when $q(z|x)$ is its own marginal; this intuitively makes for less varying conditional distributions, where $q(z|x)$ depends less on $x$. While clearly there is a trivial minima ($q(z|x)$ always the same function, completely ignoring $x$), this is tempered by a reconstruction term 


\subsection{Invariant Codes through VAE}
\label{subsec:vae}

We apply our proposed penalty to the VAE of Kingma and Welling \cite{kingma2013auto}, inspired by the similarity of the penalty in Eq.~\ref{eq:penalty} to the VAE loss function. The original VAE stems from the classical unsupervised task of constructing latent factors, $z$, so that $p(z), p(x|z)$ define a generative model that maximizes the log likelihood of the data $\mathbb{E}_x[\log p(x)]$. This generally intractable expression is lower bounded using Jensen's inequality and a variational approximation:
\begin{align}
\log p(x) & \geq  -KL[~q(z|x)~\|~p(z)~] + \mathbb{E}_{z \sim q(z|x)}[\log p(x|z)].\label{eq:ae_loss}
\end{align}
Kingma and Welling \cite{kingma2013auto} frame $q(x|z)$ and $p(z|x)$ as an encoder/decoder pair. They then provide a re-parameterization trick that, when used with standard function approximators (neural networks), allows for efficient estimation of latent codes $z$. In short, the reparametrization is the following:
\begin{align}
q(z|x) = g_{\theta}(x) + \varepsilon, ~~~~~~~~~~~~ \varepsilon \sim \mathcal{N}(0,\sigma(\theta))
\end{align}
where $g_\theta$ is a deterministic function (neural network) with parameters $\theta$, and $\varepsilon$ is an independent random variable from a Normal distribution\footnote{In the original paper, this was defined more generally; here we only consider the Normal distribution case.} also parameterized by $\theta$.



% %We start with a generative derivation of Variational Auto-Encoders \cite{kingma2013auto} (VAE).
% Let $\{x_i\}_{i=1}^N$ be a dataset generated by $\{z_i\}$ latent factors.
% \[p(x,z) = p(z)p(x|z).\]
% The classical unsupervised task is to find $p(z), p(x|z)$ that maximize $\mathbb{E}_x[\log p(x)]$. Unfortunately, this is in general intractable, due to the intractability of $p(x) = \int p(x,z) dz$. Instead, this problem is commonly reframed into a variational setting using an arbitrary function $q$ and Jensen's Inequality:
% \begin{align}
% \log p(x) &
%   = \log \int p(x,z) dz
%   = \log \int \frac{p(x,z)}{q(z|x)} q(z|x) dz
%   = \log \mathbb{E}_{z \sim q(z|x)}[ \frac{p(x,z)}{q(z|x)} ]\\
% %
% & \geq \mathbb{E}_{z \sim q(z|x)} [ \log p(x,z) - \log q(z|x) ]\\
% & = \mathbb{E}_{z \sim q(z|x)} [ \log p(z) - \log q(z|x) ] + \mathbb{E}_{z \sim q(z|x)}[\log p(x|z)]
% \end{align}



% \subsection{Equivariant Codes}
% \label{subsec:eq-codes}
We can reformulate Kingma and Welling's VAE to include our penalty on $I(z,c)$.
Define $\{x_i\}_{i=1}^N$ data and latent factors $\{z_i\}$, but also define observed $\{c_i\}$ upon which $x$ may have non-trivial dependence. That is,
\begin{align}
p(x,z,c) = p(z,c)p(x|z,c).
\end{align}
The invariant coding task is to find $q(z|x),p(x|z,c)$ that maximize $\mathbb{E}_{(x,c)}[\log p(x|c)]$, subject to $z \perp c $ under $z\sim q$ (i.e. subject to the estimated code $z$ being invariant to $c$). We make the same relaxation as in Eq. \ref{eq:gen_loss} to formulate our objective:
\begin{align}
\max \mathbb{E}_{(x,c)}[\log p(x|c)] - \lambda I(z,c).
\end{align}
Starting with the first term, we can derive a familiar looking encoder/decoder loss function that now includes $c$.
\begin{align}
\log p(x|c) & = \log \int p(x,z|c) dz\\
& = \log \int \frac{p(x,z|c)}{q(z|x)} q(z|x) dz = \log \mathbb{E}_{z\sim q}\left[ \frac{p(x,z|c)}{q(z|x)} \right]\\
& \geq \mathbb{E}_{z\sim q}[ \log p(x,z|c) - \log q(z|x) ]\\
& = \mathbb{E}_{z\sim q}[\log p(z|c) - \log q(z|x) +  \log p(x|z,c) ] .
\end{align}
Because $p(z|c)$ is a prior, we can make the assumption that $p(z|c) = p(z)$, the prior marginal distribution over $z$. This is a willful model misspecification: for an arbitrary encoder, the latent factors, $z$, are probably not independent of $c$. However, practically we wish to find $z$ that \emph{are} independent of $c$, thus it is reasonable to include such a prior belief in our generative model.
%\BlueComment{Of course, the $z$'s generated under the posterior, $q(z|x)$ might not be satisfy this constraint, but our penalty will encourage them to do so.}
Taking this assumption, we have
\begin{align}
\log p(x|c) & \geq \mathbb{E}_{z\sim q}[\log p(z) - \log q(z|x)] + \mathbb{E}_{z\sim q}[\log p(x|z,c)]\\
& = -KL[~q(z|x)~\|~p(z)] + \mathbb{E}_{z\sim q}[\log p(x|z,c)]. \label{eq:first_term}
\end{align}
This is almost exactly the same as the VAE objective in Eq. \ref{eq:ae_loss}, except our decoder $p(x|z,c)$ requires $c$ as well as $z$. Putting this together with the penalty term Eq. \ref{eq:penalty}, we have the following variational bound on the combined objective (up to a constant):
\begin{align}
\mathbb{E}_{(x,c)}[\log P(x|c)] - \lambda I(z,c) & \geq \nonumber \\
 \mathbb{E}_{(x,c)}[ ~ -KL[~  q(z|x)~\|&~p(z)]  - \lambda KL[~q(z|x)~\|~q(z)~] +
(1+\lambda)\mathbb{E}_{z\sim q}[\log p(x|z,c)] ~~ ]. \label{eq:unsup_objective}
\end{align}
We use this bound to learn $c$-invariant auto-encoders.
%\BlueComment{$1+\lambda$ could come out... then we would get $\alpha=1/(1+\lambda)$ and we would see that we get a mixture of $\alpha A + (1-\alpha) B$.}

\subsubsection{Derivation of an approximation for the Conditional-Marginal divergence}
Equation \ref{eq:unsup_objective} is our desired loss function for learning invariant codes $z$ in an unsupervised context. Unfortunately it contains $q(z)$, the empirical marginal distribution of latent code $z$, which is difficult to compute. Using the re-parameterization trick, $q(z)$ becomes a mixture distribution and this allows us to approximate its divergence from $q(z|x)$. 
% \begin{align}
% KL[ q(z|x) \| q(z) ] & = -H(q(z|x)) - \int q(z|x) \log \int q(z|x')p(x') dx' dz\\
% & \approx -H(q(z|x)) - \int q(z|x) \log \frac{1}{B}\sum_{x'} q(z|x') dz\\
% & = -H(q(z|x)) - \mathbb{E}_{z\sim q | x}[\log \sum_{x'} q(z|x') ] - \log B\\
% & \leq -H(q(z|x)) - \sum_{x'}\mathbb{E}_{z\sim q | x} [\log q(z|x')] - \log B\\
% & = \sum_{x'}[ -H(q(z|x)) - \mathbb{E}_{z\sim q | x} [\log q(z|x')] ] + (B-1)H(q(z|x)) - \log B\\
% & = \sum_{x'} \underbrace{KL[q(z|x) \| q(z|x')]}_{\text{KL between Gaussians}} + (B-1)\underbrace{H(q(z|x))}_{\text{Gaussian Ent.}} - \log B \label{eq:KLqzxqz}
% \end{align}
\begin{align}
KL[ q(z|x) \| q(z) ] & = -H(q(z|x)) - \int q(z|x) \log \int q(z|x')p(x') dx' dz\\
& = -H(q(z|x)) - \int q(z|x) \log \mathbb{E}_{x'}[q(z|x')] dz\\
& \leq -H(q(z|x)) - \mathbb{E}_{z\sim q | x}[\mathbb{E}_{x'}[\log q(z|x') ]]\\
& \approx -H(q(z|x)) - \sum_{x'}\mathbb{E}_{z\sim q | x} [\log q(z|x')]\\
& = \sum_{x'}[ -H(q(z|x)) - \mathbb{E}_{z\sim q | x} [\log q(z|x')]]\\
& \approx \sum_{x}\sum_{x'} \underbrace{KL[q(z|x) \| q(z|x')]}_{\text{KL between Gaussians}} \label{eq:KLqzxqz}
\end{align}
We can thus approximate $KL[ q(z|x) \| q(z) ]$ from its pairwise distances $KL[q(z|x) \| q(z|x')]$, which all have a closed form due to the reparameterization trick. While this requires $O(b^2)$ operations for batch size $b$, we can reduce pairwise Gaussian KL divergence to matrix algebra, making this computation fast in practice.

This further provides insight into the previously proposed Variational Fair Auto-Encoder of Louizos et al \cite{louizos2015variational}. In that paper, the authors add a Maximum Mean Discrepancy penalty as a somewhat ad hoc regularizer. This nevertheless works in practice quite well, as it encourages the statistical moments of each $q(z|c)$ to be the same over the varying values of $c$. Our condition of $KL[ q(z|x) \| q(z) ]$ has equivalent minima, and shares the ``q-regularizing'' flavor of the MMD penalty.

\subsubsection{Alternate derivation leads to adversarial loss}
\label{sec:adv_loss}

In Equation \ref{eq:fork_in_road} we used the identity $I(z,c) = I(z, x) - I(z, x | c) + I(z, c |x)$, with the caveat that the third term $I(z,c|x)$ is zero. We could have instead used another identity, $I(z, c) = H(c) - H(c|z)$. Here, the first term is constant, but expanding the second term provides the following:
\begin{align}
H(c|z) & = \mathbb{E}_{c,z\sim q}[ -\log p(c|z) ]\\
& = \inf_{r(c|z)} \mathbb{E}_{c,z\sim q}[ -\log r(c|z) ]\\
\mathbb{E}_{(x,c)}[\log P(x|c)] - \lambda I(z,c) & \geq \nonumber \\
\mathbb{E}_{(x,c)}[ ~ -KL[~  q(z|x)~\|&~p(z)] + \mathbb{E}_{z\sim q}[\log p(x|z,c)]~~] + \lambda \inf_{r(c|z)} \mathbb{E}_{c,z\sim q}[ -\log r(c|z) ] \label{eq:adv_inf}
\end{align}
The last inequality is again up to a constant term. Interpreting this in machine learning parlance, another possible approach for minimizing $I(z, c)$ is to optimize the conditional distribution $p(c|z)$ so that $r$, the lowest entropy predictor of $c$ given $z$, has the highest entropy (i.e. is as inaccurate as possible at predicting $c$). This is often operationalized by adversarial learning, and subsequent error may be due in part to the adversary not achieving the infimum. Practically speaking, this may indicate that over-training adversaries would benefit performance by bringing the adversarial gradient closer to the infimum adversary's gradient. 

\subsection{Supervised Invariant Codes through the Variational Information Bottleneck}
\label{subsec:sup}

Learned encodings are often used for downstream supervised prediction tasks. Just as in Variational Fair Auto-encoders \cite{louizos2015variational}, we can model both at the same time to offer $c$-invariant predictions. Our formulation of this problem fits into the Information Bottleneck framework \cite{tishby2000information} and mirrors the Variational Information Bottleneck (VIB) \cite{alemi2016deep}.

Conceptually, VAEs have strong connections to the Information Bottleneck \cite{alemi2016deep}. Stepping out of the generative context, we can ``reroute'' our decoder to a label variable $y$. This gives us the following computational model:
\begin{align}
x \xrightarrow{q(z|x)} z \xrightarrow{p(y|z)} y
\end{align}
The bottleneck paradigm prescribes optimizing over $q$ and $p$ so that $I(z,y)$ is maximal while minimizing $I(x,z)$ (``maintaining information about y with maximal compression of x into z''). As illustrated by Alemi et al. \cite{alemi2016deep}, this can be approximated using variational inference.

We can produce $c$-invariant codes in the supervised Information Bottleneck context using the relaxation from Eq. \ref{eq:gen_loss}. Beginning with the bottleneck objective $\max_{q,p} I(z,y) - \beta I(x,z)$ and then including the minimization of $I(z,c)$, we have
\begin{align}
\max_{p,q} I(z,y) - \beta I(x,z) - \lambda I(z,c)
\end{align}
We can then apply the same bound as in Eq. \ref{eq:breakdown} to obtain, up to constant $H(x|c)$, the following:
\begin{align}
\max_{p,q} I(z,y) - (\beta + \lambda) I(x,z) + \lambda\mathbb{E}[\log p(x|z,c) ].
\end{align}
In this objective we have a maximization of likelihood $p(x|z,c)$. This is a decoder loss, adding a third branch to our network. Following the derivation in Alemi et al. \cite{alemi2016deep} as well as a similar path as in Section \ref{subsec:vae}, the variational bound on the objective is
\begin{align}
I(z,y) - (\beta + \lambda) I(x,z) + \lambda\mathbb{E} &[\log p(x|z,c) ] \geq \nonumber\\
\mathbb{E}_{(x,c)}[ ~ \mathbb{E}_{z,y}[\log p(y|z)]  - &(\beta+\lambda) KL[~q(z|x)~\|~q(z)~] +
\lambda\mathbb{E}_{z}[\log p(x|z,c)] ~~ ]. \label{eq:sup_objective}
\end{align}
We use Eq. \ref{eq:sup_objective} to learn $c$-invariant predictors. Optimization is performed over three function approximations: one encoder $q(z|x)$, one conditional decoder $p(x|z,c)$, and one predictor $p(y|z)$. We further must compute $KL[~q(z|x)~\|~q(z)~]$ from the $I(x,z)$ penalty term. Instead of following Alemi et al.\cite{alemi2016deep}, we again use the approximation to $KL[~q(z|x)~\|~q(z)~]$ from Eq. \ref{eq:KLqzxqz}.

%\subsection{Computational Model}

